\documentclass[11pt,a4paper]{article}\n\usepackage[utf8]{inputenc}\n\usepackage[T1]{fontenc}\n\usepackage{lmodern}\n\usepackage{amsmath,amssymb}\n\usepackage{graphicx}\n\usepackage[ruled,vlined]{algorithm2e}\n\usepackage{hyperref}\n\usepackage{enumitem}\n\usepackage{caption}\n\title{Explainability Frameworks for Neural Cryptanalysis: A Multi-Layer Attribution and Causal Validation Study on Speck32/64}\n\author{}\n\date{}\n\begin{document}\n\maketitle\n\n\begin{enumerate}[leftmargin=*]\n  \item Introduction\n\end{enumerate}\n\n\subsection{Background and Motivation}\n\nin\n\nThe rapid convergence of deep learning and classical cryptanalysis has introduced a new direction the analysis of symmetric ciphers. Traditionally, cryptanalysis relied on mathematically derived techniques such as differential and linear cryptanalysis, where distinguishing characteristics were manually constructed from the cipher's design. This paradigm shifted significantly when Gohr demonstrated that a deep residual neural network could outperform the best known classical differential distinguishers on reduced-round Speck32/64 and successfully support key recovery attacks [1]. Unlike traditional methods, the neural network was provided with no explicit cryptographic knowledge and instead learned statistical patterns directly from ciphertext pairs. While this demonstrated the remarkable capability of neural distinguishers, it also raised an important question: what cryptographic structures has the model actually learned, and how are its predictions formed?\n\nUnderstanding the internal behaviour of neural distinguishers is essential for both cryptographic analysis and trustworthy artificial intelligence. High classification accuracy alone cannot determine whether a model has learned genuine cipher properties or merely exploited statistical artefacts introduced during training. Existing research has largely focused on improving distinguishing performance, extending attack rounds, and enhancing key recovery, while comparatively little attention has been given to interpreting the decision-making process of these models. This work addresses that gap by developing a layered explainability framework that combines Integrated Gradients, causal validation, interaction analysis, structural stability studies, and control experiments to systematically investigate the reasoning process of the DeepSpeck neural distinguisher. By connecting attribution-based explanations with verifiable cryptographic behaviour, the proposed framework aims to improve both the interpretability and the scientific reliability of neural cryptanalysis.\n\n\subsection{Problem Statement}\n\nThe central problem motivating this research is that neural distinguishers for lightweight block ciphers have become increasingly powerful, yet the reasoning behind their predictions remains largely opaque. While existing models are capable of distinguishing ciphertext pairs with high accuracy and, in some cases, supporting key-recovery attacks, they provide little insight into how\n\ntheir decisions are formed. It remains unclear which ciphertext bits contribute most to the prediction, whether these contributions correspond to genuine cryptographic structures or merely statistical artefacts, how different influential bits interact with one another, and whether the learned representations remain stable as the complexity of the cipher increases. This lack of interpretability limits the scientific understanding of neural cryptanalysis and raises important questions regarding the reliability and trustworthiness of these models.\n\nAttribution methods can identify input features that appear important, but they do not establish whether those features exert a genuine causal influence on the model's predictions. Similarly, high classification accuracy alone does not indicate whether a neural distinguisher has learned meaningful cryptographic properties or simply exploited patterns specific to the training data.\n\nMotivated by these limitations, this work seeks to answer a series of fundamental research questions. Like, which ciphertext bits contribute most significantly to the decision of the DeepSpeck neural distinguisher? Do the features highlighted by attribution methods genuinely influence the model's predictions when subjected to causal intervention? How closely do attribution-based importance scores align with causally validated importance measures? Do influential bits operate independently, or do they exhibit cooperative interactions that shape the model's learned representations? Are these learned patterns preserved as the number of encryption rounds increases, indicating stable cryptographic behaviour rather than round-specific artefacts? Can carefully designed control experiments confirm that the recovered attribution patterns arise from genuine cryptographic structure instead of incidental statistical correlations?\n\nAddressing these questions forms the central objective of the proposed explainability framework and establishes the foundation for the methodology presented in this work.\n\n\subsection{Existing Approaches and Their Limitations}\n\nThe most closely related work focuses on understanding what neural distinguishers actually learn. Benamira et al. demonstrated that the internal representations of Gohr's network closely approximate the differential distribution table (DDT) of Speck rather than arbitrary statistical correlations, while Gohr, Leander, and Neumann [17] reported similar findings for Simon, providing strong evidence that neural distinguishers capture genuine cryptographic structure. the theoretical understanding of neural Although cryptanalysis, they relied primarily on layer activations and statistical analysis rather than producing interpretable, per-bit explanations or validating their findings through causal intervention.\n\nthese studies significantly\n\nimproved\n\nA parallel line of research has concentrated on improving the performance of neural distinguishers. Chen et al. [2] enhanced ciphertext representations using features derived from multiple ciphertext pairs, Yue and Wu [3] proposed an Inception-inspired architecture that achieved 97.13\\% accuracy on a seven-round Speck32/64 distinguisher, and Wang et al. [4] introduced a dual-difference training strategy validated across Speck, Simon, Simeck and Hight. In the context of key recovery, Li et al. [5] extended neural-assisted attacks to 64-bit block ciphers through a Key Bit Sensitivity Test, Lu et al. [11] developed improved related-key distinguishers for Simon and Simeck, and Chen, Shen, and Yu [6] combined neural distinguishers with classical statistical techniques to broaden the applicability of neural cryptanalysis. While these contributions substantially improve distinguishing accuracy and attack capability, they provide limited insight into the reasoning process of the underlying models.\n\nRecent advances in explainable artificial intelligence provide a potential solution to this limitation. Arrieta et al. [18] established a comprehensive taxonomy of XAI methods, while Dwivedi et al. [19] identified faithfulness as the primary criterion for evaluating explanation quality. Integrated Gradients [20] and SHAP [21] have emerged as widely adopted attribution techniques, and recent surveys [22]–[24] have demonstrated their growing use in cybersecurity. However, existing studies also highlight that attribution scores alone do not guarantee causal relevance or domain-specific validity. Consequently, no prior work has applied attribution-based explainability to neural cryptanalysis in a framework that combines bit-level attribution, causal validation, interaction analysis, cross-round stability evaluation, and control experiments to verify that the recovered explanations correspond to genuine cryptographic structures rather than statistical artefacts.\n\n\subsection{Central Thesis}\n\nThis research advances the following thesis: the internal decision-making process of neural distinguishers for lightweight block ciphers can be made interpretable and verifiable through a principled multi-layer framework that combines gradient-based attribution, causal ablation, faithfulness evaluation, interaction analysis, structural stability testing, and control validation. Concretely, we argue that when Integrated Gradients attribution scores are correlated with causal ablation results and shown to be disrupted by deliberate destruction of the cipher's differential structure, the resulting explanations constitute genuine, cryptographically grounded evidence about the features the model relies on — rather than post hoc rationalisations that happen to look plausible. We further argue that tracking the evolution of these explanations across round configurations reveals whether the model is learning stable differential propagation patterns or round-specific artefacts, providing a longitudinal perspective that neither mechanistic activation analysis nor single-round attribution can supply on its own.\n\nThe framework developed here does not seek to replace classical cryptanalysis. Rather, it seeks to make neural distinguishers scientifically transparent, so that their outputs can be critically evaluated, their conclusions can be mapped back to specific structural properties of a cipher, and their failures can be understood rather than merely observed.\n\n\subsection{Contributions of Prior Work}\n\nThe present research builds on a substantial body of prior contributions across two interconnected fields. In neural cryptanalysis, Gohr [1] established the foundational architecture and input format for neural distinguishers; follow-up studies by Chen et al. [2], Yue and Wu [3], and Wang et al. [4] established that input encoding and architecture choices independently govern performance; Li et al. [5], Chen, Shen, and Yu [6], and Lu et al. [11] extended neural distinguishers into full key-recovery frameworks; and Benamira et al. and Gohr, Leander, and Neumann [17] provided mechanistic evidence that distinguishers approximate the DDT. In the XAI domain, Sundararajan et al. [20] established the axiomatic foundations of Integrated Gradients; Lundberg and Lee [21] developed SHAP as a unified game-theoretic attribution framework; Arrieta et al. [18] and Dwivedi et al. [19] established taxonomies and evaluation criteria for XAI methods; and Yan et al. [23] and Capuano et al. [24] surveyed XAI applications in cybersecurity. These bodies of work collectively provide the theoretical foundations and methodological toolkit on which the present framework draws.\n\n\subsection{Contributions of This Work}\n\nThis paper makes the following novel contributions to the field of neural cryptanalysis:\n\nMulti-layer explainability framework: We design and implement an eight-layer framework that progressively transforms the DeepSpeck neural distinguisher from a black-box classifier into an interpretable system. The layers — data generation, attribution, causal validation, faithfulness evaluation, interaction analysis, structural stability, randomised control, and structure destruction control — collectively provide a comprehensive account of the model's decision-making process.\n\nPer-bit attribution: We apply Integrated Gradients to the 64-bit input representation of Speck32/64 ciphertext pairs and produce a global importance map aggregated across 5,000 samples, revealing a highly non-uniform attribution distribution dominated by a small subset of ciphertext bits (prominently bits 24, 26, 22, 21, and 20).\n\n\begin{verbatim}\nCausal  validation  via single-bit ablation: We confirm that the attribution hotspots correspond\nto  genuinely  influential  features  through  systematic  single-bit  ablation  experiments,  in  which\n\end{verbatim}\n\nremoving individual bits produces accuracy drops of up to 8.44 percentage points from a baseline of approximately 64\\%.\n\nPairwise interaction: We examine higher-order feature dependencies through pairwise ablation experiments, constructing an interaction matrix that reveals synergistic and redundant bit pairs. The strongest synergies are observed between bit pairs (22, 18), (22, 54), and (18, 54), suggesting that the distinguisher exploits distributed, multi-bit cryptographic structures.\n\nStructural stability analysis across rounds: We compare attribution maps generated from models trained on five, six, seven, and eight rounds of Speck32/64, computing round-to-round correlation scores and tracking the transport of attribution motifs. The analysis reveals that while exact hotspot locations shift between rounds, the overall attribution structure is preserved with increasing stability at higher round counts — evidence that the model learns persistent differential propagation patterns.\n\nControl validation: We conduct two control experiments — randomised structure validation and true structure destruction — which confirm that the recovered attribution patterns are specific to genuine cryptographic structure and do not reflect dataset artefacts. Destruction of differential (49.78\\%) while producing structure fundamentally different attribution maps, validating the framework end-to-end.\n\nreduces classification accuracy\n\n\begin{verbatim}\nto  near-random\n\end{verbatim}\n\n\begin{enumerate}[leftmargin=*]\n  \item Related Work\n\end{enumerate}\n\nThe pivotal starting point is Gohr's 2019 paper [1], which demonstrated that a deep residual network trained on ciphertext pairs generated by Speck32/64 could achieve distinguishing accuracy substantially beyond classical differential distinguishers, reaching 5–8 effective rounds and enabling 11–12-round key recovery. Gohr's model accepted a concatenated ciphertext pair as input and was trained end-to-end without explicit domain knowledge, marking a fundamental departure from algebraic cryptanalysis. This architecture and input format—raw final-round ciphertext pairs, no auxiliary information—serves as the baseline configuration for the present project.\n\nA consistent theme across follow-up work is that both network architecture and input representation independently govern distinguisher accuracy. Chen et al. [2], published in The Computer Journal (2023), introduced an input format aggregating derived features from multiple ciphertext pairs rather than a single pair, significantly improving accuracy on Speck32/64 and Simon32/64. A study in Applied Sciences [3] replaced Gohr's initial convolutional block with an Inception-inspired residual block, reporting 97.13\\% accuracy for a 7-round Speck32/64 distinguisher—though part of this gain likely stems from an enriched input format including\n\nearlier-round difference information rather than the architectural change alone. Wang et al. [4], in IET Information Security (2024), proposed feeding ciphertext pairs generated under two distinct input differences simultaneously, with an automated four-step pipeline (difference selection, sample generation, training, evaluation) validated on Simon, Speck, Simeck, and Hight. Collectively, these results establish that comparisons of raw accuracy across papers must account for differences in input encoding—a methodological consideration directly relevant to interpreting attribution results in the present project.\n\nSeveral journal papers have extended neural distinguishers into full attack pipelines. Li et al. [5], in Frontiers of Information Technology and Electronic Engineering (2024), introduced a Key Bit Sensitivity Test to partition the key space for large-state block ciphers, extending key recovery from toy 32-bit ciphers to operationally relevant 64-bit variants. Chen, Shen, and Yu [6], in The Computer Journal (2023), combined neural distinguishers with classical statistical techniques for a neural-aided statistical attack, demonstrating that the paradigm extends beyond pure differential cryptanalysis. Lu et al. [11], also in The Computer Journal (2024), produced improved related-key differential neural distinguishers for Simon and Simeck competitive with classical search-based approaches. Together, these works confirm that neural distinguishers are important components of practical key-recovery frameworks—making it that much more important for understanding which features drive their decisions.\n\nNeural distinguishers have also been used as an evaluation tool for cipher designs. Tian and Hu [14], in KSII Transactions on Internet and Information Systems (2021), applied deep-learning-assisted differential cryptanalysis to Simon, showing that ciphers with provable diffusion properties remain susceptible to neural distinguishers trained without structural priors. The Journal of Information Security and Applications (2023) carried a study evaluating IoT-targeted ciphers LBC-IoT and SLIM against neural distinguishers [15], and a related study by Liu et al. [16] proposed a multi-round, multi-splicing construction that systematically improves distinguisher accuracy by chaining multiple cipher rounds as input. These works illustrate that neural distinguishers are now treated as a standard cryptanalytic benchmark, reinforcing the need for interpretability tools that can verify whether their conclusions reflect genuine cipher weaknesses.\n\nA critical gap in this literature is the absence of principled understanding of what these models have actually learned. Benamira et al., in an analysis of Gohr's distinguisher, probed intermediate activations and showed that the model learns to approximate the differential distribution table (DDT) rather than arbitrary statistical correlations. Gohr, Leander, and Neumann's assessment of differential-neural distinguishers [17] similarly concluded that the\n\nneural distinguisher for Simon functions as a close approximation of the DDT, providing a theoretical baseline against which XAI-recovered feature importance scores can be evaluated.\n\nThese mechanistic analyses confirm that neural distinguishers encode genuine cryptographic structure in a distributed, implicit form. However, they do not provide per-bit attribution scores explaining which ciphertext bit positions drive classification, nor how this changes as round count increases. This is precisely the gap the present project addresses by applying Integrated Gradients, SHAP, and LIME to trained neural distinguishers, and validating the results causally.\n\nExplainable AI (XAI) has matured significantly as a research discipline. Arrieta et al. [18], in Information Fusion (2020), published a comprehensive taxonomy classifying XAI approaches into ante hoc and post hoc categories, identifying feature attribution as the dominant paradigm for explaining deep neural networks in high-stakes domains. Dwivedi et al. [19], in ACM Computing Surveys (2023), established that faithfulness—the degree to which an explanation reflects the model's actual computation—is the central criterion for evaluating attribution quality.\n\nIntegrated Gradients (IG) [20] computes attributions by tracing how the model's output changes as the input is gradually moved from a baseline (e.g., all-zero ciphertext) to the actual input, with each bit's score reflecting its share of that change. Unlike raw gradient methods, IG guarantees that a bit which genuinely affects the prediction always receives a non-zero score, and that the scores sum exactly to the model's output difference. In the cryptographic setting, this gives each ciphertext bit a score reflecting its contribution to the distinguisher's decision.\n\nA critical review in Artificial Intelligence Review (2026) [22] analyzed faithfulness, stability, and operational applicability across mainstream XAI techniques, concluding that no single method dominates on all criteria and that combining multiple attribution methods—as proposed here—provides stronger evidence of genuine feature importance than any single approach. Yan et al. [23], in the International Journal of Intelligent Systems (2022), and Capuano et al. [24], in IEEE Access (2022), both surveyed explainable machine learning in cybersecurity contexts, cataloguing applications of attribution-based XAI methods to intrusion detection, malware classification, and encrypted traffic analysis. Both identify a shared limitation directly relevant to this project: attribution explanations may identify features that are statistically important to a classifier without being genuinely meaningful from the domain perspective—exactly the faithfulness problem this project addresses through causal interventions.\n\nNeural distinguishers are now well-established cryptanalytic tools, with documented efficacy across Speck, Simon, Simeck, Hight, and other lightweight ciphers, and mechanistic analyses confirm they encode genuine cryptographic structure approximating the DDT. However, no existing work has applied attribution-based XAI methods directly to neural distinguishers to produce per-bit importance scores, validate those scores causally, or study how explanations evolve across round counts and cipher families. The XAI literature demonstrates that Integrated Gradients and SHAP are suitable for binary feature attribution in security applications, and that causal-intervention-based faithfulness validation is an established evaluation paradigm. The proposed framework applies these techniques to neural cryptanalysis for the first time, bridging the performance-focused neural cryptanalysis literature and the interpretability-focused XAI literature.\n\n\begin{enumerate}[leftmargin=*]\n  \item Proposed Methodology\n\end{enumerate}\n\nThis research is designed to uncover, validate, explain and interpret the internal decision processes of a black box model, here being neural cryptanalysis models—especially those designed to differentiate between random data and cipher outputs. The methodology is bifurcated into three stages, each one of them aiming to identify crucial input features, ensure their meaningfulness and certainty in the context of cipher.\n\n\begin{enumerate}[leftmargin=*]\n  \item Explainability via Integrated Gradients (IG): We set about by applying integrated\n\end{enumerate}\n\ngradients to our model. Integrated gradients put forth an upright measure of each input bit’s contribution to the model’s output. Using this method, we obtain a detailed ranking of which bits were the most influential in decision making by the model. This is done on a large number of samples, say 5000, making sure that the attribution is stable, significant and representative of the model’s global behaviour.\n\n\begin{enumerate}[leftmargin=*]\n  \item Causal Validation through Single-Bit Ablation: After having a list of important bits,\n\end{enumerate}\n\nwe understand their influence on the model using causal ablation. This phase systematically masks every important bit, one at a time, and measures the overall drop in model accuracy. Higher the accuracy drop, crucial is the bit. We compare the significance among the significant bits we produced in the previous phase, thus filtering out the “false positives”.\n\n\begin{enumerate}[leftmargin=*]\n  \item Multi-Bit Interaction and Structural Analysis: In our final phase, we explore the\n\end{enumerate}\n\nrelationship between the top-ranked bits. We understand their dependence, correlation and interaction that can change the decision of the model. Bit pairs are examined and their combined accuracy is measured. If the joint ablation of two bits causes a greater\n\nimpact on accuracy than the sum of their individual drops, then we discover a collaborative interaction between the two bits. Apart from this, we test if our model withholds its observations and discoveries with increased complexity. This helps us understand whether the neural model captures stable structures or is just a mere illusion.\n\n\begin{verbatim}\nF ( x ) + F ( y ) < F ( x , y )\n\end{verbatim}\n\nWhere F(x) is the ablation function of x.\n\nBy following the above methodology, we validate the model’s decision with justifiable explanations and proofs that suit them. The entire research is based on different case studies that not only answer numerous questions but also ensures the robustness of the model across cryptographic configurations.\n\n\subsection{Speck32/64 Cipher Architecture}\n\nThe foundation of this work is the Speck32/64 lightweight block cipher. Speck was developed by the United States National Security Agency (NSA) as a family of lightweight cryptographic algorithms intended for resource-constrained environments. Unlike conventional block ciphers that rely on complex substitution boxes and lookup tables, Speck was specifically designed to achieve high security while maintaining extremely low computational and memory requirements.\n\nThe variant used throughout this study is Speck32/64, which operates on a 32-bit plaintext block and uses a 64-bit encryption key. The plaintext block is divided into two 16-bit words. During encryption, these words undergo a sequence of repeated transformations known as rounds. Each round gradually transforms the original plaintext into ciphertext while preserving reversibility through the use of carefully designed mathematical operations. Speck has embedded ARX structure in it, where ARX refers to Addition, Rotation, and XOR. This structure is more efficient on modern processors as it creates highly complex dependencies between bits as encryption progresses.\n\n\subsection{Neural Cryptanalysis of Speck model}\n\nConventionally, differential cryptanalysis relies on manual derivation of differential characteristics to bifurcate ciphertext from random data. But recent discoveries in deep learning have shown that neural networks can learn these patterns using large datasets automatically. We have used the publicly available DeepSpeck neural distinguisher model introduced by Gohr in our research. The neural distinguisher is designed to determine whether a given ciphertext pair originates from the speck encryption process or from random data.\n\nThe model is formulated as a binary classification task. Ciphertext pairs generated using valid differential characteristics are assigned a positive label, while randomly generated pairs are assigned a negative label. The neural network learns to identify subtle statistical structures that distinguish the two classes. The classifier outputs a probability score ranging between 0 and 1. Values close to 1 indicate that the input pair is likely generated by Speck, whereas values close to 0 indicate random data.\n\nFormally, the model learns a function:\n\n\begin{verbatim}\nF(X)→[0,1]\n\end{verbatim}\n\nwhere X represents the binary representation of the ciphertext pair and the output represents the confidence of the classifier.\n\n\subsection{Input Representation}\n\nThe neural network converts the integer ciphertext values to binary representations before supplying it to the model. The Speck32/64 operates on two 16-bit words and two ciphertexts are used per sample.\n\n4 × 16 = 64 binary features.\n\n\begin{verbatim}\nTherefore, every sample is represented as a 64-dimensional binary vector.\n\end{verbatim}\n\n\subsection{Requirements and Specifications}\n\nThe model was trained using differential ciphertext pairs generated from the Speck cipher. Rather than relying on handcrafted features, the network automatically learns discriminative representations through supervised learning.\n\nThe specific model analyzed in this work is: net7\\_small.h5\n\n\begin{verbatim}\nwhich corresponds to a neural distinguisher trained on 7-round Speck32/64 data.\n\end{verbatim}\n\nThe model serves as the foundation for all explainability experiments conducted in this study. Several software libraries have been put in use to implement the proposed framework. Such as tensor flow, Keras, Numpy, Matplotlib and scipy.\n\n\subsection{Architecture and Framework}\n\nAny explainable model typically terminates at the attribution layer, after finding the important bits. A few of them understand their causal interpretations and validate them with visual representations. But verifying these bits and their significance is a behavioural gap that our frameworks attempt to cover. To address this limitation we developed a multi-layer explainability framework.The neural distinguisher model, introduced by Gohr, forms the base of the research, on which we apply multiple layers of case studies, interpretations and validations.\n\nThe architecture consists of the following layers:\n\n\begin{enumerate}[leftmargin=*]\n  \item Data Generation Layer\n\end{enumerate}\n\n\begin{enumerate}[leftmargin=*]\n  \item Neural Attribution Layer\n\end{enumerate}\n\n\begin{enumerate}[leftmargin=*]\n  \item Causal Validation Layer\n\end{enumerate}\n\n\begin{enumerate}[leftmargin=*]\n  \item Faithfulness Evaluation Layer\n\end{enumerate}\n\n\begin{enumerate}[leftmargin=*]\n  \item Interaction Analysis Layer\n\end{enumerate}\n\n\begin{enumerate}[leftmargin=*]\n  \item Structural Transport Layer\n\end{enumerate}\n\n\begin{enumerate}[leftmargin=*]\n  \item Randomized Control Layer\n\end{enumerate}\n\n\begin{enumerate}[leftmargin=*]\n  \item Structure Destruction Control Layer\n\end{enumerate}\n\nEvery layer contributes to the conclusions derived from the framework using perturbation, ablation, transportation, destruction, randomization, causal learning and validation. The complete framework transforms a black box model into an interpretable system.\n\n\subsubsection{Dataset Construction and Generation Layer}\n\\begin{algorithm}[H]
\\caption{Dataset Construction and Generation}
\\label{alg:data-generation}
\\KwData{Speck32/64 cipher, key generation process, input differential structure}
\\KwResult{Structured and random ciphertext pairs in 64-bit binary form}
Generate random plaintext pairs\;
Introduce the predefined input differential to construct structured samples\;
Encrypt structured and random plaintext pairs using Speck32/64\;
Convert each ciphertext pair into a 64-dimensional binary feature vector\;
Assign structured vectors to the structured class and random vectors to the random class\;
Use the generated dataset for DeepSpeck explainability experiments\;
\\end{algorithm}\n\nUnlike most of the research works that are based on publicly available fixed datasets, cryptographic data has a unique property of generating itself whenever needed, as long as the encryption algorithm, key generation process and input differentials are known. Here we use the Speck32/64 encryption algorithm to generate ciphertext pairs.\n\nSince the framework is based on binary classification, the algorithm produces two classes of ciphertext pairs. The first class contains structured samples (generated using encrypted plain texts), and the second class contains random samples. The Speck32/64 encrypts random plaintext pairs that are generated using a predefined input differential structure. The differential characteristics are preserved among the structured samples, while the random data loses their features.\n\nOnce generated, every ciphertext pair is converted into a 64-dimensional binary feature vector. Since each sample consists of four 16-bit ciphertext words, the binary representation allows the neural network to analyze relationships directly at the bit level. This helps the ARX based operations to work smoothly during the encryption process as the complexity increases.\n\nThe generated binary vectors are then classified into structured or random classes using the pretrained DeepSpeck neural distinguisher. The same dataset used throughout the process, in every layer and experiments conducted. This ensures that any observations made are due to the explainability layers and not the dataset itself.\n\nRandom Plaintext Pairs\n\n\begin{verbatim}\n          ↓\n\end{verbatim}\n\nIntroduce Differential (Structured Samples)\n\n\begin{verbatim}\n          ↓\n\end{verbatim}\n\nSpeck32/64 Encryption\n\n\begin{verbatim}\n          ↓\n\end{verbatim}\n\nCiphertext Pairs\n\n\begin{verbatim}\n          ↓\n\end{verbatim}\n\nBinary Conversion\n\n\begin{verbatim}\n          ↓\n\end{verbatim}\n\n\begin{verbatim}\n         64-Dimensional Feature Vector\n\end{verbatim}\n\n\begin{verbatim}\n          ↓\n\end{verbatim}\n\nDeepSpeck Neural Distinguisher\n\n\subsubsection{Attribution Layer}\n\\begin{algorithm}[H]
\\caption{Neural Attribution via Integrated Gradients}
\\label{alg:attribution}
\\KwData{Pretrained DeepSpeck model, ciphertext pairs dataset}
\\KwResult{Global importance map for 64 input bits}
For each ciphertext sample in the dataset\;
  Pass the sample through the pretrained DeepSpeck model\;
  Compute Integrated Gradients from a baseline input to the actual sample\;
  Record attribution scores for all 64 input bits\;
Aggregate attribution scores across samples\;
Compute the global importance map from the aggregated scores\;
\\end{algorithm}\n\nThe second layer of the framework aims to answer the most fundamental explainability question:\n\nWhich input bits influence the neural network's decisions?\n\nTo answer this question, Integrated Gradients (IG) was employed. Integrated Gradients is a gradient-based attribution technique that measures how much each input feature contributes to a model's prediction. Unlike ordinary gradients, which often suffer from instability and saturation effects, Integrated Gradients accumulate gradients along a continuous path between a baseline input and the actual sample.\n\nFor each ciphertext sample:\n\n\begin{enumerate}[leftmargin=*]\n  \item A dataset of ciphertext pairs was generated using the Speck32/64 cipher.\n\end{enumerate}\n\n\begin{enumerate}[leftmargin=*]\n  \item Samples were passed through the pretrained DeepSpeck model.\n\end{enumerate}\n\n\begin{enumerate}[leftmargin=*]\n  \item Integrated Gradients was computed for each sample.\n\end{enumerate}\n\n\begin{enumerate}[leftmargin=*]\n  \item Attribution scores were calculated for all 64 input bits.\n\end{enumerate}\n\n\begin{enumerate}[leftmargin=*]\n  \item Scores were aggregated across multiple samples.\n\end{enumerate}\n\n\begin{enumerate}[leftmargin=*]\n  \item A global importance map was generated.\n\end{enumerate}\n\nThis process is repeated across thousands of samples.\n\nThe attribution scores are then averaged to obtain a global importance map.\n\nThe resulting importance distribution reveals which regions of the ciphertext the neural network considers most relevant during classification.\n\nThis layer produces the first explainability artifact of the framework: the hotspot map.\n\n\subsubsection{Causal Validation Layer}\n\\begin{algorithm}[H]
\\caption{Causal Validation by Single-Bit Ablation}
\\label{alg:causal-validation}
\\KwData{Pretrained DeepSpeck model, attribution-ranked input bits}
\\KwResult{Causal importance scores for each input bit}
Measure the baseline accuracy of the neural distinguisher\;
For each input bit\;
  Remove or mask the bit in the input samples\;
  Re-evaluate the modified samples with the DeepSpeck model\;
  Compute the change in classification accuracy\;
  Record the accuracy loss as the causal importance score\;
\\end{algorithm}\n\nThe previous layer tells us the significant bits but doesn't give reasons that can justify their importance. It doesn’t prove the impact of these bits on the model’s decision.Therefore, the third layer introduces causal validation.\n\nWhat happens when the model is deprived of its most important input bits?\n\nThe objective of this stage is to test whether the importance discovered by Integrated Gradients corresponds to genuine influence on model behavior.The procedure is straightforward.\n\nFor each input bit:\n\n\begin{enumerate}[leftmargin=*]\n  \item The baseline accuracy of the neural distinguisher was recorded.\n\end{enumerate}\n\n\begin{enumerate}[leftmargin=*]\n  \item Each input bit was removed individually.\n\end{enumerate}\n\n\begin{enumerate}[leftmargin=*]\n  \item The modified samples were re-evaluated by the model.\n\end{enumerate}\n\n\begin{enumerate}[leftmargin=*]\n  \item The change in classification accuracy was measured.\n\end{enumerate}\n\n\begin{enumerate}[leftmargin=*]\n  \item Accuracy loss was recorded as the causal importance score.\n\end{enumerate}\n\nIf removing a bit causes a substantial reduction in accuracy, that bit is considered causally important.\n\nIf the accuracy remains unchanged, the bit is considered non-essential regardless of its attribution score.\n\n\subsubsection{Faithfulness Verification Layer}\n\\begin{algorithm}[H]
\\caption{Faithfulness Verification of Explanations}
\\label{alg:faithfulness}
\\KwData{Attribution scores, causal importance scores}
\\KwResult{Faithfulness measures for the explanation}
Compare attribution scores and causal importance scores for all input bits\;
Calculate a statistical correlation coefficient between both distributions\;
Compute top-K overlap between the highest-ranked attribution bits and the highest-ranked causal bits\;
Assess faithfulness using correlation and overlap metrics\;
\\end{algorithm}\n\nCan the explanations given by the model be trusted? Do their fidelity scores justify their accuracy? Do the explanations generated by the model faithfully reflect its actual decision-making process?\n\nTo answer the above questions, this layer was introduced. Once attribution scores and causal scores are obtained, the next challenge is determining whether they agree.\n\nThis layer evaluates the faithfulness of the explanations.\n\nFaithfulness refers to the degree to which an explanation reflects the true behavior of the model.\n\nA faithful explanation should satisfy the following condition:\n\nHigh Attribution\n\n\begin{verbatim}\n↓\n\end{verbatim}\n\nHigh Causal Impact\n\nTo evaluate this relationship, the framework compares:\n\n\begin{itemize}[leftmargin=*]\n  \item Integrated Gradient scores\n\end{itemize}\n\n\begin{itemize}[leftmargin=*]\n  \item Causal ablation scores\n\end{itemize}\n\nStatistical correlation analysis is then performed between both distributions.\n\nAdditionally, overlap analysis is conducted between the most important attribution bits and the most important causal bits. Quantitative approach like top-K overlap is used to measure the extent to which the features highlighted by the attribution method correspond to those that genuinely influence the model's predictions, thereby providing an indication of the faithfulness of the generated explanations.\n\nThe faithfulness of explanations is decided using correlation coefficients, statistical significance value and top-K overlap scores. This stage is the first quantitative assessment of explanation quality.\n\n\subsubsection{Interaction Analysis Layer}\n\\begin{algorithm}[H]
\\caption{Interaction Analysis of Important Bit Pairs}
\\label{alg:interaction-analysis}
\\KwData{Top-ranked important bits, pretrained DeepSpeck model}
\\KwResult{Interaction matrix for bit pairs}
Select pairs of the highest-ranked bits\;
For each selected bit pair\;
  Remove both bits simultaneously from the input samples\;
  Evaluate the modified samples with the DeepSpeck model\;
  Record the resulting classification accuracy\;
  Compare the combined accuracy loss to the sum of individual bit losses\;
Aggregate the pairwise interaction scores into an interaction matrix\;
\\end{algorithm}\n\nCryptographic algorithms rarely depend on a single bit acting independently. Through the repeated mixing steps used in encryption, information from many bits becomes tightly intertwined as it passes through successive rounds. As a result, a neural distinguisher may uncover relationships that arise from multiple bits working together, rather than from any one bit alone. Because of this, evaluating bits in isolation can only provide a partial picture of how the model makes its decisions. This subsection therefore explores that issue in detail.\n\nDo the most important bits operate independently, or do they work together to\n\ninfluence the model's decisions?\n\nTo investigate these hidden dependencies, an Interaction Analysis Layer was introduced after the causal validation stage. The primary goal of this layer is to determine whether the neural network relies on combinations of important bits and to identify how these bits interact with one another during classification.\n\nThe following steps are performed within this layer:\n\n\begin{enumerate}[leftmargin=*]\n  \item The highest-ranked bits are selected pairwise\n\end{enumerate}\n\n\begin{enumerate}[leftmargin=*]\n  \item For each possible pair of important bits, both bits are simultaneously removed from the\n\end{enumerate}\n\ninput data\n\n\begin{enumerate}[leftmargin=*]\n  \item The modified samples are then passed through the neural distinguisher\n\end{enumerate}\n\n\begin{enumerate}[leftmargin=*]\n  \item The resulting classification accuracy is recorded\n\end{enumerate}\n\n\begin{enumerate}[leftmargin=*]\n  \item This combined accuracy loss is compared against the accuracy drops observed when the\n\end{enumerate}\n\nsame bits were removed individually\n\n\begin{enumerate}[leftmargin=*]\n  \item The interaction scores obtained from all bit pairs are aggregated into an interaction matrix\n\end{enumerate}\n\nthat captures the dependency structure learned by the model\n\nBy comparing these values, it becomes possible to measure whether the effect of removing two bits together is greater than, smaller than, or approximately equal to the sum of their individual effects.\n\nIf removing both bits causes a substantially larger performance degradation than expected, the pair is considered to exhibit a cooperative or synergistic relationship. This indicates that the neural network relies on information shared between those bits. Conversely, if the combined effect is smaller than expected, the bits may contain overlapping information, suggesting redundancy within the learned representation.\n\n\subsubsection{Structural Stability Layer}\n\\begin{algorithm}[H]
\\caption{Structural Stability Across Encryption Rounds}
\\label{alg:structural-stability}
\\KwData{Neural distinguishers trained on different round configurations}
\\KwResult{Correlation of importance maps across rounds}
For each round configuration model\;
  Generate the importance map for its ciphertext input bits\;
Extract the highest-ranked bits from each importance map\;
Compute correlation between importance distributions of different round configurations\;
Compare early, intermediate, and later rounds\;
Analyze how important regions shift, expand, contract, or migrate across rounds\;
\\end{algorithm}\n\nDo the learned cryptographic patterns remain stable as the number of encryption\n\nrounds increases?\n\nAfter identifying the single bit, their correlations and synergistic structures, it is important to understand their longevity. Can these discovered patterns survive complex ciphers?\n\n\begin{itemize}[leftmargin=*]\n  \item If the same regions of the ciphertext continue to influence the model across multiple\n\end{itemize}\n\nencryption rounds, it suggests that the neural network has learned genuine cryptographic structures.\n\n● Conversely, if the important regions change completely from one round configuration to another, the model may be relying on round-specific artifacts rather than fundamental properties of the cipher.\n\nTo investigate this, a Structural Stability Layer was introduced within the framework. The purpose of this layer is to track how the model's attention changes as the number of encryption rounds increases and to determine whether important patterns persist throughout the encryption process.\n\nThe analysis is performed in the following stages:\n\n\begin{enumerate}[leftmargin=*]\n  \item Multiple neural distinguishers trained on different Speck round configurations are\n\end{enumerate}\n\nselected\n\n\begin{enumerate}[leftmargin=*]\n  \item Every model is analyzed separately with a different level of cryptographic complexity\n\end{enumerate}\n\n\begin{enumerate}[leftmargin=*]\n  \item Every round configuration is performed independently and respect importance maps are\n\end{enumerate}\n\ngenerated\n\n\begin{enumerate}[leftmargin=*]\n  \item Each map represents how strongly the model relies on every ciphertext bit while making\n\end{enumerate}\n\nits predictions\n\n\begin{enumerate}[leftmargin=*]\n  \item The highest-ranked bits from each attribution map are extracted and analyzed\n\end{enumerate}\n\n\begin{enumerate}[leftmargin=*]\n  \item Correlation analysis is performed between the importance distributions obtained from\n\end{enumerate}\n\ndifferent round configurations\n\n\begin{enumerate}[leftmargin=*]\n  \item The comparison is carried out between:\n\end{enumerate}\n\n\begin{itemize}[leftmargin=*]\n  \item Early and intermediate rounds\n\end{itemize}\n\n\begin{itemize}[leftmargin=*]\n  \item Intermediate and later rounds\n\end{itemize}\n\n\begin{itemize}[leftmargin=*]\n  \item Early and later rounds\n\end{itemize}\n\n\begin{enumerate}[leftmargin=*]\n  \item The framework then examines how important regions evolve across rounds\n\end{enumerate}\n\n\begin{enumerate}[leftmargin=*]\n  \item We study how groups of important bits shift, expand, contract, or migrate between\n\end{enumerate}\n\nneighboring positions, leading to a phenomenon known as pattern transport\n\n\begin{enumerate}[leftmargin=*]\n  \item The resulting correlation values are used to assess structural stability.\n\end{enumerate}\n\n● High correlation indicates that the model continues to focus on similar regions across multiple rounds. This suggests that the neural distinguisher has learned persistent cryptographic structures rather than isolated patterns.\n\n● Moderate correlation indicates partial stability, where some important regions remain relevant while others evolve as the cipher becomes more complex.\n\n● Low correlation indicates that the model changes its attention significantly between round configurations. Such behavior may suggest reliance on round-specific features or unstable decision-making strategies.\n\nThe Structural Stability Layer therefore provides a longitudinal view of the neural network's behavior. Rather than analyzing importance at a single point in time, it examines how learned representations evolve throughout the encryption process. This makes it possible to determine whether the explanations discovered by the framework represent stable cryptographic knowledge or temporary patterns that emerge only under specific round configurations.\n\n\subsubsection{Control Validation Layer}\n\\begin{algorithm}[H]
\\caption{Control Validation Experiments}
\\label{alg:control-validation}
\\KwData{Original and modified ciphertext datasets, pretrained DeepSpeck model}
\\KwResult{Validation of explanation specificity}
\\textbf{Randomized Control Validation:}\;
Randomize the underlying structure of the input data while preserving its format\;
Run the explainability pipeline on randomized data\;
Compare hotspot patterns and attribution maps to the original dataset\;
Measure the similarity and divergence of importance distributions\;
\\textbf{True Structure Destruction Control:}\;
Destroy the cryptographic structure while maintaining the same statistical dimensions\;
Repeat the explainability pipeline on the destroyed-structure dataset\;
Compare the resulting attribution maps to the original dataset\;
Assess whether hotspot patterns disappear or change significantly\;
\\end{algorithm}\n\nThe final stage of the framework consists of control experiments. We need to know that our explanations were based on relevant patterns and structures. Control experiments are essential because they test whether the discovered explanations remain valid under extreme conditions.\n\nWithout controls, it is impossible to determine whether observed patterns represent genuine cryptographic structure or random artifacts.\n\nFor this reason, two independent control experiments were performed.\n\n\begin{itemize}[leftmargin=*]\n  \item Randomized Control Validation\n\end{itemize}\n\nIn this experiment, the underlying structure of the input data is randomized while preserving the overall data format.\n\nThe purpose is to evaluate whether attribution hotspots emerge even when meaningful cryptographic information is absent.\n\nIf similar hotspot patterns appear after randomization, the explanations may be unreliable.\n\nConversely, substantial differences indicate that the discovered hotspots are likely linked to genuine cryptographic structure.\n\nThis experiment acts as a sanity check for the entire explainability pipeline.\n\n\begin{itemize}[leftmargin=*]\n  \item True Structure Destruction Control\n\end{itemize}\n\nThis experiment represents the most rigorous validation stage of the framework.\n\nInstead of randomizing selected components, the cryptographic structure itself is intentionally destroyed.\n\nThe modified dataset retains the same dimensions and statistical format but no longer preserves the meaningful relationships generated by the Speck encryption process.\n\nThe explainability pipeline is then repeated.\n\nIf the discovered hotspots disappear or change significantly, it provides strong evidence that the original explanations were linked to real cipher structure.\n\nIf the hotspots remain unchanged, the explanations may be capturing dataset artifacts rather than cryptographic behavior.\n\nThis experiment provides the strongest form of validation within the framework.\n\n\subsubsection{Framework Summary}\n\nThe proposed framework progressively transforms neural cryptanalysis from a prediction task into an interpretable scientific investigation.The complete workflow can be summarized as:\n\nData Generation\n\n\begin{verbatim}\n          ↓\n\end{verbatim}\n\nIntegrated Gradients Attribution\n\n\begin{verbatim}\n          ↓\n\end{verbatim}\n\nCausal Validation\n\n\begin{verbatim}\n          ↓\n\end{verbatim}\n\nFaithfulness Verification\n\n\begin{verbatim}\n          ↓\n\end{verbatim}\n\nInteraction Analysis\n\n\begin{verbatim}\n          ↓\n\end{verbatim}\n\nStructural Stability Analysis\n\n\begin{verbatim}\n          ↓\n\end{verbatim}\n\nRandomized Control Validation\n\n\begin{verbatim}\n          ↓\n\end{verbatim}\n\nStructure Destruction Validation\n\n\begin{verbatim}\n↓\n\end{verbatim}\n\nStatistical Interpretation\n\nEach layer contributes a distinct form of evidence. Together, these layers provide a comprehensive evaluation of whether the neural distinguisher has learned genuine cryptographic structure and whether the generated explanations can be trusted.\n\n3.6. Testing methodology\n\n\subsubsection{Experimental Setup}\n\nThe proposed explainability framework was evaluated using the publicly available DeepSpeck neural distinguisher for the Speck32/64 block cipher. The objective of the experimental setup was to provide a controlled environment for analyzing the internal decision-making process of the neural cryptanalysis model.\n\nThe experimental setup consisted of the following components:\n\n\begin{enumerate}[leftmargin=*]\n  \item Neural Model\n\end{enumerate}\n\n\begin{itemize}[leftmargin=*]\n  \item Pretrained DeepSpeck neural distinguisher\n\end{itemize}\n\n\begin{itemize}[leftmargin=*]\n  \item Trained for ciphertext pair classification\n\end{itemize}\n\n\begin{itemize}[leftmargin=*]\n  \item Used as the base model for all explainability experiments\n\end{itemize}\n\n\begin{enumerate}[leftmargin=*]\n  \item Programming Environment\n\end{enumerate}\n\n\begin{itemize}[leftmargin=*]\n  \item Python\n\end{itemize}\n\n\begin{itemize}[leftmargin=*]\n  \item TensorFlow / Keras\n\end{itemize}\n\n\begin{itemize}[leftmargin=*]\n  \item NumPy\n\end{itemize}\n\n\begin{itemize}[leftmargin=*]\n  \item Statistical analysis libraries\n\end{itemize}\n\nThese tools were used for model execution, attribution analysis, causal interventions, and result evaluation.\n\n\begin{enumerate}[leftmargin=*]\n  \item Dataset Generation\n\end{enumerate}\n\n\begin{itemize}[leftmargin=*]\n  \item Dataset generated directly using the Speck32/64 encryption algorithm\n\end{itemize}\n\n\begin{itemize}[leftmargin=*]\n  \item Structured ciphertext pairs generated using differential cryptanalysis principles\n\end{itemize}\n\n\begin{itemize}[leftmargin=*]\n  \item Random ciphertext pairs generated as the negative class\n\end{itemize}\n\n\begin{itemize}[leftmargin=*]\n  \item Each sample converted into a 64-bit binary feature vector\n\end{itemize}\n\n\begin{enumerate}[leftmargin=*]\n  \item Experimental Consistency\n\end{enumerate}\n\n\begin{itemize}[leftmargin=*]\n  \item The same data generation procedure was maintained throughout the framework\n\end{itemize}\n\n\begin{itemize}[leftmargin=*]\n  \item All explainability experiments were performed on a controlled input distribution\n\end{itemize}\n\n\begin{itemize}[leftmargin=*]\n  \item This ensured fair comparison between different analysis layers\n\end{itemize}\n\n5. Evaluation Strategy The framework was not evaluated solely on prediction accuracy. Instead, multiple aspects of model behavior were examined:\n\n\begin{itemize}[leftmargin=*]\n  \item Bit-level importance\n\end{itemize}\n\n\begin{itemize}[leftmargin=*]\n  \item Causal influence\n\end{itemize}\n\n\begin{itemize}[leftmargin=*]\n  \item Explanation faithfulness\n\end{itemize}\n\n\begin{itemize}[leftmargin=*]\n  \item Feature interactions\n\end{itemize}\n\n\begin{itemize}[leftmargin=*]\n  \item Structural stability across rounds\n\end{itemize}\n\n\begin{itemize}[leftmargin=*]\n  \item Robustness under control experiments\n\end{itemize}\n\nBy combining a benchmark neural cryptanalysis model with controlled data generation and multiple validation procedures, the experimental setup provides a reliable environment for investigating how and why the neural distinguisher makes its predictions.\n\n\subsubsection{Attribution Evaluation Methodology}\n\nHow was feature importance identified within the neural distinguisher?\n\nThe Attribution Layer was evaluated using the Integrated Gradients (IG) technique. The objective of this experiment was to identify which ciphertext bits contributed most strongly to the model's predictions. The results that validated this layer were:\n\n\begin{itemize}[leftmargin=*]\n  \item Bit-level importance scores\n\end{itemize}\n\n\begin{itemize}[leftmargin=*]\n  \item Ranked list of influential bits\n\end{itemize}\n\n\begin{itemize}[leftmargin=*]\n  \item Global attribution map\n\end{itemize}\n\nThis experiment provided the initial explanation of the model's behavior and served as the foundation for all subsequent validation stages.\n\n\subsubsection{Causal Validation Methodology}\n\nAre the bits identified by Integrated Gradients genuinely responsible for the model's\n\ndecisions?\n\nTo validate the importance scores obtained from Integrated Gradients, a causal intervention experiment was performed using single-bit ablation. Certain factors were responsible for the conclusions drawn from this layer. Like,\n\n\begin{itemize}[leftmargin=*]\n  \item Accuracy drop for every input bit\n\end{itemize}\n\n\begin{itemize}[leftmargin=*]\n  \item Causal importance ranking\n\end{itemize}\n\n\begin{itemize}[leftmargin=*]\n  \item Identification of truly influential features\n\end{itemize}\n\nThis experiment transformed attribution-based observations into experimentally verified causal evidence.\n\n\subsubsection{Faithfulness Evaluation Methodology}\n\nDo the explanations generated by the model faithfully reflect its actual\n\n\begin{verbatim}\ndecision-making process?\n\end{verbatim}\n\nThe purpose of this experiment was to determine whether the attribution scores produced by Integrated Gradients aligned with the causal importance measurements obtained through ablation.\n\nProcedure:\n\n\begin{itemize}[leftmargin=*]\n  \item Integrated Gradients scores were collected.\n\end{itemize}\n\n\begin{itemize}[leftmargin=*]\n  \item Causal importance scores were collected.\n\end{itemize}\n\n\begin{itemize}[leftmargin=*]\n  \item Correlation analysis was performed between the two rankings.\n\end{itemize}\n\n\begin{itemize}[leftmargin=*]\n  \item Top-k overlap analysis was conducted.\n\end{itemize}\n\n\begin{itemize}[leftmargin=*]\n  \item Agreement scores were calculated.\n\end{itemize}\n\nEvaluation Measures:\n\n\begin{itemize}[leftmargin=*]\n  \item Pearson Correlation\n\end{itemize}\n\n\begin{itemize}[leftmargin=*]\n  \item Top-10 Overlap\n\end{itemize}\n\n\begin{itemize}[leftmargin=*]\n  \item Top-20 Overlap\n\end{itemize}\n\n\begin{itemize}[leftmargin=*]\n  \item Agreement Score\n\end{itemize}\n\nHigher agreement indicates that the generated explanations accurately reflect the model's behavior.\n\n\subsubsection{Interaction Analysis Methodology}\n\nThis experiment investigated whether the neural distinguisher relied on isolated features or groups of interacting features. Certain conclusions were drawn using the generated measures below,\n\n\begin{itemize}[leftmargin=*]\n  \item Pairwise interaction matrix\n\end{itemize}\n\n\begin{itemize}[leftmargin=*]\n  \item Synergy scores\n\end{itemize}\n\n\begin{itemize}[leftmargin=*]\n  \item Redundancy scores\n\end{itemize}\n\nThe experiment revealed hidden dependencies among important features learned by the model.\n\n\subsubsection{Structural Stability Evaluation Methodology}\n\nThis experiment examined whether the patterns identified by the model persisted across different levels of encryption complexity. Comparisons performed were evaluated using,\n\n\begin{itemize}[leftmargin=*]\n  \item Round-to-round correlation scores\n\end{itemize}\n\n\begin{itemize}[leftmargin=*]\n  \item Pattern transport measurements\n\end{itemize}\n\n\begin{itemize}[leftmargin=*]\n  \item Stability indicators\n\end{itemize}\n\nThis experiment assessed whether the model had learned persistent cryptographic structures or round-specific patterns.\n\n\subsubsection{Control Validation Methodology}\n\nControl experiments were conducted using randomized data. Below given are the steps it followed to conduct this experiment.\n\n\begin{itemize}[leftmargin=*]\n  \item Attribution maps were generated using genuine ciphertext samples.\n\end{itemize}\n\n\begin{itemize}[leftmargin=*]\n  \item A second dataset with randomized structure was created.\n\end{itemize}\n\n\begin{itemize}[leftmargin=*]\n  \item The same explainability procedures were applied to both datasets.\n\end{itemize}\n\n\begin{itemize}[leftmargin=*]\n  \item Importance distributions were compared.\n\end{itemize}\n\n\begin{itemize}[leftmargin=*]\n  \item Entropy and similarity measures were calculated.\n\end{itemize}\n\nEvaluation Measures:\n\n\begin{itemize}[leftmargin=*]\n  \item Entropy\n\end{itemize}\n\n\begin{itemize}[leftmargin=*]\n  \item Correlation\n\end{itemize}\n\n\begin{itemize}[leftmargin=*]\n  \item Importance Distribution Similarity\n\end{itemize}\n\nThis experiment served as a sanity check and verified whether the framework was identifying meaningful cryptographic patterns rather than random noise.\n\n\begin{enumerate}[leftmargin=*]\n  \item RESULTS AND OUTCOMES\n\end{enumerate}\n\n\subsection{Layerwise Explainability Results}\n\nIn this section, results, tables, graphs and comparison discussion of respective layers are attached. The objective was to evaluate the model using the proposed explainability framework and identify the cryptographic structures learned during classification.\n\n\subsubsection{Attribution Layer Results}\n\nIntegrated Gradients was applied to the pretrained DeepSpeck model to identify the relative importance of individual input bits. Attribution scores were aggregated across multiple samples to generate a global importance map.\n\nTable 4.1: Rankwise Bit Positions and IG scores\n\nFigure 4.1: Integrated Gradients Heatmap\n\nFigure 4.2: Integrated Gradients Variance Analysis\n\nDiscussion\n\n\begin{verbatim}\nUnexpectedly, the attribution analysis revealed a highly non-uniform distribution of\n\end{verbatim}\n\nimportance across the 64 input bits. A small subset of bits dominated the attribution map, with Bits 24, 26, 22, 21, and 20 receiving the highest Integrated Gradients scores. The most influential bit achieved a normalized importance score of 1.0, while the remaining top-ranked bits also exhibited substantially higher contributions than the average feature.\n\nThe heatmap visualizations clearly show the presence of concentrated attribution hotspots rather than evenly distributed importance. Furthermore, the variance analysis indicates that these hotspots remain consistent across multiple samples, suggesting that the neural distinguisher repeatedly relies on the same ciphertext regions during classification. This concentration of importance provides initial evidence that the model has learned structured cryptographic patterns rather than random statistical correlations.\n\n\subsubsection{Causal Validation Results}\n\nSingle-bit ablation experiments were performed to verify whether the bits identified by attribution analysis were causally important. Each bit was removed individually and the resulting change in model performance was measured.\n\nTable 4.2: Top Causally Important Bits\n\nFigure 4.3: Bitwise Causal Importance\n\nFigure 4.4: Hotspot Ablation Accuracy\n\nFigure 4.5: Hotspot Ablation Confidence\n\nDiscussion\n\nThe causal validation results demonstrate that several attribution hotspots identified by\n\nIntegrated Gradients correspond to genuinely influential features. The model achieved a baseline accuracy of 64.26\\%, which decreased significantly when certain bits were removed. In particular, Bits 22, 54, 36, and 4 produced the largest accuracy drops, with Bit 22 alone reducing performance by 8.44\\%.\n\nSeveral of these bits also appeared among the highest-ranked attribution features, indicating consistency between attribution and intervention-based analysis. Furthermore, the confidence measurements show a noticeable reduction in prediction certainty after removing important bits, confirming that the neural distinguisher actively relies on these regions during classification rather than merely assigning them high attribution scores.\n\n\subsubsection{Faithfulness Verification Results}\n\nThe attribution rankings obtained from Integrated Gradients were compared against the causal rankings obtained through ablation analysis.\n\nTable 4.3: Evaluation Metrics\n\nFigure 4.6: Integrated Gradients vs Causal Importance\n\nFigure 4.7: Attribution–Causal Correlation Analysis\n\nDiscussion\n\nA positive correlation was observed between the attribution scores and causal importance values, indicating that the features identified by Integrated Gradients generally correspond to bits that genuinely influence model performance. The attribution–causal correlation coefficient of 0.265, along with Top-10 and Top-20 overlaps of 4/10 and 7/20 respectively, demonstrates a measurable level of agreement between the two methods. Notably, bits such as 22, 54, 18, and 50 appeared among the highest-ranked features in both analyses. While the agreement is not perfect, the results suggest that the generated explanations capture a significant portion of the model's decision-making process and provide reasonably faithful interpretations of the neural distinguisher's behavior.\n\n\subsubsection{Interaction Analysis Results}\n\nTo investigate higher-order dependencies, pairwise ablation experiments were performed on the most important bits identified during causal validation.\n\nFigure 4.8: Pairwise Interaction Heatmap\n\nFigure 4.9: Top Feature Interactions\n\nDiscussion\n\nThe interaction analysis revealed the presence of both synergistic and redundant feature\n\npairs within the neural distinguisher. Out of the evaluated interactions, 23 pairs exhibited negative synergy while only 4 pairs showed positive synergy, indicating that the model primarily relies on distributed and interdependent feature groups.\n\nThe strongest interactions were observed between bit pairs (22,18), (22,54), and (18,54), each producing substantially larger performance degradation when removed together than when removed individually. These results suggest that the distinguisher does not base its predictions solely on isolated ciphertext bits but instead exploits relationships among multiple features, reflecting the complex dependency structures introduced by the Speck encryption process.\n\n\subsubsection{Structural Stability Results}\n\nAttribution maps generated from models trained on different round configurations were compared to evaluate the persistence of learned patterns.\n\nTable 4.4: Dominant Bits Per Round\n\n\begin{verbatim}\nTable 4.5: Round-to-Round Correlation\n\end{verbatim}\n\nFigure 4.10: Motif Transport Heatmap\n\nFigure 4.11: Motif Transport Correlation Analysis\n\nFigure 4.12: Motif Transport Structure Visualization\n\n\begin{verbatim}\nFigure 4.13: Round-wise Comparison\n\end{verbatim}\n\nDiscussion\n\n\begin{verbatim}\nThe round-wise attribution maps exhibit a measurable degree of structural persistence\n\end{verbatim}\n\nacross increasing encryption depths. Correlation values indicate that a significant portion of the learned importance structure is retained as additional rounds are introduced. The highest correlation observed between the 8-round and 9-round models suggests that the neural distinguisher increasingly stabilizes its focus on specific cryptographic regions at higher complexities.\n\nThe motif transport and round comparison visualizations further show that while the exact hotspot locations may shift across rounds, the overall attribution structure remains preserved. This behavior indicates that the model is tracking persistent cryptographic motifs rather than learning isolated round-specific patterns. The results therefore suggest that the neural distinguisher captures stable features associated with differential propagation within the Speck cipher, providing evidence that the learned representations reflect genuine cipher behavior rather than transient artifacts of a particular round configuration.\n\n\subsubsection{Control Validation Results}\n\nControl experiments were conducted by comparing attribution patterns obtained from genuine ciphertext data against randomized inputs.\n\nFigure 4.14: Real vs Random Comparison\n\nFigure 4.15: Real vs Random Variance Analysis\n\nFigure 4.16: Real vs Randomized Heatmap\n\nFigure 4.17: Real vs Randomized Structure Analysis\n\nDiscussion\n\nThe control experiments revealed clear differences between genuine ciphertext data and randomized inputs. The entropy increased from 16.994 to 17.799 after randomization, indicating that attribution became more dispersed when cryptographic structure was removed. Although a high correlation of 0.968 was observed between the two distributions, several hotspot rankings changed, suggesting that the model relies on meaningful cryptographic patterns rather than purely random features. These results support the validity of the proposed explainability framework and demonstrate its ability to capture genuine cipher-specific behavior.\n\n\subsection{Case Study 1 - Differential Structure Removal}\n\nThe objective of this case study is to determine whether the DeepSpeck neural distinguisher relies on the differential characteristics embedded within the ciphertext pairs. To evaluate this, the original Speck-generated ciphertext pairs were replaced with randomly generated binary inputs while keeping the model unchanged. If the neural distinguisher has learned meaningful cryptographic structure, its performance should deteriorate significantly after the differential information is removed.Two datasets were evaluated:\n\n\begin{itemize}[leftmargin=*]\n  \item Original Dataset: Ciphertext pairs generated using the standard DeepSpeck differential\n\end{itemize}\n\ndata generation process.\n\n\begin{itemize}[leftmargin=*]\n  \item Randomized Dataset: Random binary vectors having the same dimensions as the\n\end{itemize}\n\noriginal ciphertext pairs but containing no differential structure.\n\nThe pretrained DeepSpeck neural distinguisher was tested on both datasets, and the resulting accuracy and feature distributions were compared.\n\nTable 4.6: Results of Differential Structure Removal Experiment\n\nFigure 4.18: Effect of Differential Structure Removal on Classification Accuracy\n\nFigure 4.19: Variance Distribution Across Input Bits\n\nFigure 4.20: Top Important Bit Comparison\n\nDiscussion\n\nThe results demonstrate a clear reduction in model performance after the differential structure was removed. As shown in Figure 4.18, the classification accuracy decreased from 61.56\\% on the original dataset to 49.78\\% on the randomized dataset. Since 50\\% accuracy corresponds to random guessing in a binary classification task, this indicates that the neural distinguisher loses almost all predictive capability when the underlying differential relationships are destroyed.\n\nFigure 4.19 shows the variance distribution across all 64 input bits. The two curves exhibit nearly identical variance levels, with mean variances of 0.249956 and 0.249958 respectively. This observation is important because it confirms that the accuracy degradation is not caused by changes in the statistical distribution of the inputs. Both datasets possess similar variance characteristics, yet only the structured ciphertext pairs allow successful classification.\n\nFigure 4.20 compares the highest-ranked bit positions obtained from the two datasets. The important bit locations identified in the original dataset differ considerably from those observed after randomization. The absence of a stable hotspot structure in the randomized case suggests that the meaningful feature patterns learned by the model are directly associated with the differential properties of the Speck cipher. Collectively, these results indicate that the neural distinguisher is not exploiting arbitrary statistical correlations. Instead, it derives its predictive capability from genuine cryptographic structures present within the differential ciphertext pairs.\n\n\subsection{Case Study 2 - Noise Injection Analysis}\n\nThe objective of this case study is to evaluate the robustness of the DeepSpeck neural\n\ndistinguisher when the input ciphertext pairs are corrupted with random noise. Unlike the\n\nprevious case study, where the differential structure was completely removed, this experiment preserves the underlying cryptographic structure while introducing controlled perturbations. The goal is to determine whether the learned features remain reliable when a portion of the input information is altered.\n\nA noise rate of 5\\% was applied to the input dataset. For each ciphertext pair, individual bits were randomly selected and flipped with a probability of 0.05. The pretrained DeepSpeck neural distinguisher was then evaluated on both the original and noisy datasets.\n\nThe following metrics were recorded:\n\n\begin{itemize}[leftmargin=*]\n  \item Classification Accuracy\n  \item Prediction Confidence\n\end{itemize}\n\nFigure 4.23: Impact of Noise Injection on Classification Accuracy\n\nFigure 4.24: Prediction Confidence Under Noise\n\nTable 4.7: Noise Injection Results\n\nDiscussion\n\nThe introduction of 5\\% random noise resulted in a noticeable reduction in model performance. As shown in Figure 4.23, the classification accuracy decreased from 62.86\\% to 54.02\\%, corresponding to an accuracy drop of 8.84 percentage points. This indicates that even a relatively small amount of corruption can significantly affect the neural distinguisher's ability to recognize differential patterns.\n\nA similar trend can be observed in Figure 4.24, where the mean prediction confidence decreased from 0.1156 to 0.1077 after noise injection. The reduction in confidence suggests that the model becomes less certain about its predictions when the input structure is partially distorted.\n\nDespite the presence of noise, the model still performs above random guessing, indicating that a substantial portion of the learned cryptographic information remains intact. However, the observed performance degradation demonstrates that the distinguisher relies on stable bit-level structures that are sensitive to perturbations. These results highlight the importance of preserving cryptographic feature integrity for reliable neural cryptanalysis.\n\nReferences\n\n\begin{verbatim}\n[1]  A. Gohr, "Improving attacks on round-reduced Speck32/64 using deep learning," in\n\end{verbatim}\n\nAdvances in Cryptology – CRYPTO 2019, vol. 11693, pp. 150–179, Springer, 2019.\n\n[2] Y. Chen, Y. Shen, H. Yu, and S. Yuan, "A new neural distinguisher considering features derived from multiple ciphertext pairs," The Computer Journal, vol. 66, no. 6, pp. 1419–1433, 2023. https://doi.org/10.1093/comjnl/bxac019\n\n[3] Yue, X., \\& Wu, W. (2023). Improved Neural Differential Distinguisher Model for\n\nLightweight Cipher Speck. Applied Sciences, 13(12), 6994. https://doi.org/10.3390/app13126994\n\n\begin{verbatim}\n[4]  G. Wang et al., "A new (related-key) neural distinguisher using two differences for\n\end{verbatim}\n\ndifferential cryptanalysis," IET Information Security, 2024. https://doi.org/10.1049/2024/4097586\n\n[5] X. Li, J. Ren, and S. Chen, "Improved deep learning aided key recovery framework:\n\napplications to large-state block ciphers," Frontiers of Information Technology and Electronic Engineering, vol. 25, no. 10, pp. 1406–1420, 2024. https://doi.org/10.1631/FITEE.2300848\n\n\begin{verbatim}\n[6]  Y. Chen, Y. Shen, and H. Yu, "Neural-aided statistical attack for cryptanalysis," The\n\end{verbatim}\n\nComputer Journal, vol. 66, no. 10, pp. 2480–2498, 2023.https://doi.org/10.1093/comjnl/bxac099\n\n\begin{verbatim}\n[11] J. Lu, G. Liu, B. Sun, et al., "Improved (related-key) differential-based neural\n\end{verbatim}\n\ndistinguishers for SIMON and SIMECK block ciphers," The Computer Journal, vol. 67, no. 2, pp. 537–547, 2024.\n\n[14] W. Tian and B. Hu, "Deep learning assisted differential cryptanalysis for the lightweight cipher Simon," KSII Transactions on Internet and Information Systems, vol. 15, no. 2, pp. 600–616, 2021. https://doi.org/10.3837/tiis.2021.02.012\n\n[15] Teng, W.J., Teh, J.S., \\& Jamil, N. (2023). On the security of lightweight block\n\nciphers against neural distinguishers: Observations on LBC-IoT and SLIM. J. Inf. Secur. Appl., 76, 103531.\n\n[16] J. Liu, J. Ren, S. Chen, and M. Li, "Improved neural distinguishers with\n\nmulti-round and multi-splicing construction," Journal of Information Security and Applications, vol. 74, p. 103461, 2023.\n\n\begin{verbatim}\n[17]  A. Gohr, G. Leander, and P. Neumann, "An assessment of differential-neural\n\end{verbatim}\n\ndistinguishers," IACR Cryptology ePrint Archive, Report 2022/1521, 2022.\n\n[18] A. B. Arrieta et al., "Explainable Artificial Intelligence (XAI): Concepts, taxonomies, opportunities, and challenges toward responsible AI," Information Fusion, vol. 58, pp. 82–115, 2020. https://doi.org/10.1016/j.inffus.2019.12.012\n\n[19] R. Dwivedi et al., "Explainable AI (XAI): Core ideas, techniques, and solutions,"\n\nACM Computing Surveys, vol. 55, no. 9, pp. 1–33, 2023.\n\n[20] M. Sundararajan, A. Taly, and Q. Yan, "Axiomatic attribution for deep networks," in\n\nProceedings of ICML, PMLR, 2017. (Foundational methods paper; original IG axioms have no journal-form alternative.)\n\n\begin{verbatim}\n[21]  S. M. Lundberg and S.-I. Lee, "A unified approach to interpreting model\n\end{verbatim}\n\npredictions," in\n\nAdvances in NeurIPS, vol. 30, 2017. (Foundational methods paper for SHAP.)\n\n\begin{verbatim}\n[22]  Lin, T., Lau, R.Y.K. & Hu, S. A critical review of state-of-the-art explainable\n\end{verbatim}\n\nartificial intelligence (XAI) methods and their business applications. Artif Intell Rev 59, 161 (2026). https://doi.org/10.1007/s10462-026-11565-y\n\n[23] J. Yan et al., "Explainable machine learning in cybersecurity: A survey," International Journal of Intelligent Systems, vol. 37, no. 12, pp. 10744–10803, 2022. https://doi.org/10.1002/int.23088\n\n[24] N. Capuano et al., "Explainable artificial intelligence in cybersecurity: A survey,"\n\nIEEE Access, vol. 10, 2022. https://doi.org/10.1109/ACCESS.2022.3204171\n\n\end{document}